\chapter{Hygiene}
\label{chp:HYG}
\gAasRsN{RS~II}

% Hyg-Fortbildung 2009/11 Graz

\emph{Maintainer: Sebastian Gabriel}
\marginpar{\includegraphics[width=4cm]{./images/gabriel-christine-aass/hygiene-hello-002-edited.jpg}}

\minitoc

\nocite{Longmore2007}
\nocite{IntroClinEmergMed4}
\nocite{Boecker2}
\nocite{Fauci2008}
\nocite{ForthPharma8}

% \gChangelog{
% Changelog:
% 
% \begin{description}
%     \item[2009-02-20] Fertig für Kurs 2009/02.
%     \item[2009-02-25] Start Changelog.
%     \item[2009-04-02] Übernahme aus Openoffice in \LaTeX
% \end{description}
% }

\section{Grundlagen}

\subsection{Wichtige Begriffe}

%\begin{multicols}{2}
\begin{description}
    \item[\gT{Infektion}] \index{Infektion}Eindringen von
    Krankheitserregern in den Organismus.
    \item[\gT{Pathogenität}]
    \index{Pathogenität}Fähigkeit eines Krankheitserregers, eine Krankheit
    auszulösen.
    \item[\gT{Resistenz}]
    \index{Resistenz}Widerstandsfähigkeit eines Organismus gegen die 
    krankmachenden Eigenschaften eines Erregers.
    \item[\gT{Inkubationszeit}]
    \index{Inkubationszeit}symptomlose Zeit zwischen Infektion (Eintritt d. Erregers in den Organismus) und dem Auftreten
    erster Symptome.
    \item[\gT{Epidemie}] \index{Epidemie}örtlich
    begrenztes, zeitlich vermehrtes Auftreten einer  Infektionserkrankung
    (Grippe, Salmonellen).
    \item[\gT{Pandemie}] \index{Pandemie}Epidemie die sich
    über Länder u. Kontinente ausbreitet (Pest im Mittelalter).
    \item[\gZ{Endemie}] \index{Endemie}Auf kleinere Gebiete
    beschränktes, zeitlich jedoch nicht begrenztes Auftreten einer Infektionserkrankung  (FSME,
    Malaria, Lepra).
    \item[\gT{Morbidität}] \index{Morbidität}Anteil der
    an XY \gEs{Erkrankten} einer bestimmten \gE{Bevölkerungsgruppe}
    (z.B. 5 pro 100\,000 Einwohner).
    \item[\gT{Mortalität}] \index{Mortalität}Anteil der
    an XY \gEs{Verstorbenen} einer bestimmten \gE{Bevölkerungsgruppe}
    (z.B. 5 pro 100\,000 Einwohner).
    \item[\gT{Letalität}]  \index{Letalität}Anteil der
    an XY \gEs{Verstorbenen} bezogen auf die \gE{Gesamtzahl der
    Erkrankten} ("`1 von 5 Erkrankten stirbt"').
    \item[\gT{Entzündung}] \gSee{topic:entzuendung}: \begin{inparaenum}
      \item Rötung,
      \item Schwell\-ung,
      \item Überwärmung,
      \item Schmerzen,
      \item Funktionsverlust.
    \end{inparaenum}
    \item[\gT{Prophylaxe}] \index{Prophylaxe} Vorbeugende
    Massnahme(n).
\end{description}
%\end{multicols}





%%%%%%%%%%%%%%%%%%%%%%%%%%%%%%%%%%%%%%%%%%%%%%%%%%%%%%%%%%%
%%%%%%%%%%%%%%%%%%%%%%%%%%%%%%%%%%%%%%%%%%%%%%%%%%%%%%%%%%%
%%%%%%%%%%%%%%%%%%%%%%%%%%%%%%%%%%%%%%%%%%%%%%%%%%%%%%%%%%%
\subsection{Die wichtigsten Erreger}
\label{topic:infektion-erreger}

\begin{table}[H]
\gCaptionof{table}{Übersicht der Erregerarten.}

\begin{tabulary}{\linewidth}{LL}
\gTabHeading{Erreger}
&
\gTabHeading{Kurzbeschreibung}
\\\midrule\addlinespace
\gTabSub{Bakterien}
&
klein, im Mikroskop meist sichtbar;
        {\textquotedblleft}Normale{\textquotedblright} Zellen
\\\addlinespace
\gTabSub{Viren}
&
sehr klein, im Mikroskop i.d.R. nicht
        sichtbar, keine echten Zellen, müssen Zellen befallen
\\\addlinespace
\gTabSub{Pilze}
&

\\\addlinespace
\gTabSub{Parasiten, Tiere}
&
Würmer, Maden, \ldots
\\\addlinespace
\gTabSub{(Toxine)}
&
Giftstoffe, werden z.B. von Bakterien produziert.
\\\addlinespace\bottomrule
\end{tabulary}
\end{table}

% \subsubsection{"`Biologische Arbeitsstoffe"'}
% 
% % Risikogruppen nach VbA (Verordnung biologische Arbeitsstoffe)
% %
% % 2: Influenzaviren
% % 3: SARS, Affenpockenvirus, HIV
% % 4: Ebola, Lassa, Marburg
% 
% \begin{table}[H]
% \gCaptionof{table}{Biologische Arbeitsstoffe -- Risikogruppen}
% 
% \begin{tabulary}{\linewidth}{LL}
% \gTabHeading{Gruppe}
% &
% \gTabHeading{Bedeutung}
% \\\midrule\addlinespace
% \gTabSub{1}
% &
% 
% \\\addlinespace
% \gTabSub{2}
% &
% 
% \\\addlinespace
% \gTabSub{3}
% &
% 
% \\\addlinespace
% \gTabSub{4}
% &
% 
% \\\addlinespace\bottomrule
% \end{tabulary}
% {\footnotesize nach \S\,40 Abs. 4 Z 1 -- 4 AschG}
% \end{table}

\subsubsection{Bakterien}

\index{Bakterien}

Bakterien sind vollständige,
{\quotedblbase}echte{\textquotedblleft} Zellen, die leben
und aktiv etwas tun können. Dazu gehört z.B. die Produktion
von Substanzen (Toxine ...) und Gasen, etc. 

Bakterien sind nicht grundsätzlich "`böse"'. Viele können
für den Menschen gut, egal oder manchmal auch böse sein.
Dies hängt von verschiedenen Faktoren ab, s.u.
\gZ{Bakterien haben die unterschiedlichsten Formen (Kugeln,
Stäbchen, mit Geißeln \ldots}


\gLayout{22}{Bakterien und Pathogenität}
{
Die Schädlichkeit von Bakterien hängt von verschiedenen
Faktoren ab:
    \begin{itemize}
        \item Welches Bakterium: Spezies, Stamm, Resistenzen
        \item Ort der Infektion.
        \item "`Konkurrenzkampf der Bakterien"': Welche anderen Mikroorganismen siedeln noch an
       dem Ort?
        \item Immunstatus des Patienten
        \begin{itemize}
        \item Kind, Chemotherapie, AIDS, ... 
        \item Körperliche Verfassung
        \item Alter, Ernährung, Stress, ...
        \end{itemize}
    \end{itemize}
}{}
{images/AASdev20090228-img96.jpg}
{Für manche Bakterien zahlt man
    freiwillig viel Geld, hier am Beispiel des
    Lactobacillus casei gezeigt.}
{}



\gFazit{Daraus folgt:
    \begin{itemize}
        \item Manche Keime machen immer krank.
        \item Manche Keime nur unter Umständen krank.
        \item \gEs{Opportunisten} machen nur krank wenn
        der Wirt geschwächt ist. Das sind viele
        Patienten im Rettungs- und Krankentransport!
    \end{itemize}
    }

\paragraph{Bakterien, Antibiotika und Resistenzen}

Antibiotika sind Medikamente, die Bakterien abtöten
oder deren Vermehrung verhindern. Antibiotika haben ein bestimmtes Spektrum, dh. wirken nur
auf bestimmte Keime.

Ein großes Problem ist die zunehmende Bildung von
\gEs{Resistenzen} der Bakterien gegen Antibiotika, welche
die Medikamente wirkungslos machen. Man spricht auch von der
Bildung von \gE{multiresistenten Keimen}. Diese werden
unter \gSee{topic:multiresistente-keime} näher besprochen.




\subsubsection{Viren}
\index{Viren}

Viren sind sehr klein \gZ{und im
Lichtmikroskop nicht sichtbar}. Sie sind keine vollständigen
Zellen \gZ{und brauchen Wirtszellen um sich vermehren zu
können. Sie schleusen ihr eigenes genetisches Programm
(Erbinformationen) in die Wirtszelle ein, diese produziert
dann neue Viren.}

Virale Infektionen sind sehr häufig, die meisten sog.
\gEs{"`Verkühlungskrankheiten"'} und andere bekannte
Krankheiten wie z.B.: \gEs{Grippe}, \gE{Feuchtblattern} und
viele \gE{Hepatitis}-Formen fallen darunter.




\paragraph{Nachweis} \gZ{Viren sind nicht auf Nährböden
anzüchtbar. Man verwendet entweder Antigen- oder Antikörpernachweis
    oder die PCR (DNA/RNA-Nachweis).} Die Tests haben
    üblicherweise ein \gE{"`diagnostisches Fenster"'}, in
    dem eine Infektion (noch) nicht nachweisbar ist.


%%%%%%%%%%%%%%%%%%%%%%%%%%%%%%%%%%%%%%%%%%%%%%%%%%%%%%%%%%%
%%%%%%%%%%%%%%%%%%%%%%%%%%%%%%%%%%%%%%%%%%%%%%%%%%%%%%%%%%%
%%%%%%%%%%%%%%%%%%%%%%%%%%%%%%%%%%%%%%%%%%%%%%%%%%%%%%%%%%%
\subsection{Die wichtigsten Übertragungswege}
\label{topic:infektion-uebertragungswege}

Krankheitserreger können auf unterschiedliche Wege in den
Körper gelangen und eine Infektion verursachen:

\begin{multicols}{2}
\begin{itemize}
    \item Fäkal-oral (Schmierinfektion, z.B. Badesee)
    \item Aerogen (Tröpfcheninfektion, Aerosol)
    \item Genital (Geschlechtsverkehr)
    \item Über/durch die Haut
    \item Diaplazentär (während Schwangerschaft)
    \item Während der Geburt
    \item Eintrittspforten: alles was irgendwie offen ist
    \begin{itemize}
        \item Wunden, Magen-Darm-Trakt, Atmung, Auge, Penis/Scheide, Plazenta
        \ldots
    \end{itemize}
    \item Fremdkörper sind Wohnhäuser und Autobahnen für Keime
    \begin{itemize}
        \item Harnkatheter, venöse Zugänge, Tubus, künstlicher
        Mageneingang, Dialysekatheter {\dots}
    \end{itemize}
\end{itemize}
\end{multicols}




%%%%%%%%%%%%%%%%%%%%%%%%%%%%%%%%%%%%%%%%%%%%%%%%%%%%%%%%%%%
%%%%%%%%%%%%%%%%%%%%%%%%%%%%%%%%%%%%%%%%%%%%%%%%%%%%%%%%%%%
%%%%%%%%%%%%%%%%%%%%%%%%%%%%%%%%%%%%%%%%%%%%%%%%%%%%%%%%%%%
\subsection{Kreuzinfektion}
\index{Kreuzinfektion}
\label{topic:kreuzinfektion}

\gDefinition{\gTag{Kreuzinfektion}Verschleppung von Erregern von einem Patienten
auf einen (unbeteiligten) Anderen, wobei das Personal oder
Gegenstände als "`Vehikel"' fungieren.}

\gFazit{Kreuzinfektion: \gEs{Patient 1
${\rightarrow}$ \emph{Sanitäter} ${\rightarrow}$ Patient
2}}




% \section{Wichtige Krankheitsbilder}
%\marginlabel{Krankheitsbilder}

\section{Multiresistente Keime}
\label{topic:multiresistente-keime}

\marginlabel{Linktip: \emph{EUREGIO MRSA-net}
\url{http://www.mrsa-net.org/}\cite{mrsanet}}
\begin{quote}
\gZ{Zu Beginn des 20. Jahrhunderts waren
bakterielle Infektionen sehr gefürchtet, viele Leute starben
daran. Wer es im Ersten Weltkrieg in ein Lazarett geschafft
hatte, verstarb oft dort an den Folgen des Wundbrandes.
Erkrankungen wie Scharlach oder "`Blutvergiftungen"' waren
oft ein Todesurteil.

\marginpar{\includegraphics[width=\linewidth]{./images/hygiene/Faroe_stamp_079_europe_fleming.jpg}\par
\gCaptionof{figure}{\gZ{Alexander Fleming. \gSourcePicture{PD, Postverk Føroya -
Philatelic Office}}}}Dann ließ, im Jahre 1928, Alexander Fleming eine Nährplatte
mit Bakterienkulturen versehentlich verschimmeln und
bemerkte, dass in der Umgebung des Pilzes die Bakterien
nicht wachsen wollten: das \gE{Penicillin} war geboren. 

In weiterer Folge wurden weitere Antibiotika entwickelt und
man dachte, man hätte die bakteriellen Infektionen besiegt.
Doch man irrte sich gewaltig \ldots}
\end{quote}

\gLayout{0}{Resistenzen}
{
Bakterien können Resistenzen
gegen Antibiotika entwickeln, die Antibiotika helfen dann nicht mehr!
}
{
\begin{itemize}
   \item Bakterien können Resistenzen entwickeln!
   \item Antibiotika helfen dann nicht mehr!
\end{itemize}
}
{}
{}
{}


\gFazit{Resistenzen sind ein enorm großes Problem in der
heutigen Medizin!

\ldots und damit auch im Rettungs- und Krankentransport.}

Vor der Entdeckung der Antibiotika waren bakterielle Infektionen sehr
gefürchtet, viele Menschen sind an -- nach heutiger Sicht simplen --
Infektionen gestorben. 

Heute kann man zwar solche Infektionen mit Antibiotika behandeln, jedoch
können Bakterien \gEs{Resistenzen} gegen Antibiotika
entwickeln. \gMargin{Antibiotika werden wirkungslos}Die sonst wirksamen Medikamente
wirken dann nicht mehr oder nur noch sehr abgeschwächt. Als behandelnder Arzt steht man dann hilflos vor dem
Patienten. 

\gFazit{Resistenzen machen Antibiotika wirkungslos.}

\gFazit{Bald sind wir wieder dort, wo wir 1927, vor der
Entdeckung des Penicillins, waren: Praktisch machtlos gegen
bakterielle Infektionen.}

Es muss daher das Ziel aller sein die weitere Ausbreitung
von resistenten Keimen zu verhindern. Dem Rettungs- und Krankentransport kommt dabei eine
\gE{Schlüsselrolle} zu, da es hier leicht zu
\gEs{Kreuzinfektionen} kommen kann.


\gFazit{Resistenzen sind ein enorm großes Problem in der
heutigen Medizin!

\ldots und damit auch im Rettungs- und Krankentransport, da es
hier leicht zu Kreuzinfektionen kommen kann.}

Gesundheitseinrichtungen wie Spitäler sind eine ideale
Brutstätte für Resistenzen: Es werden viele Antibiotika
verwendet, viele Patienten liegen räumlich nahe beieinander
\ldots In Spitälern finden daher gehäuft resistente Keime. 

\gLayout{2}{Im Krankentransport gibt es Problemkeime}
{
Da im Krankentransport
natürlich häufig Patienten aus Gesundheitseinrichtungen (oder zwischen
Gesundheitseinrichtungen) transportiert werden, ist man
auch hier mit dem Problem der resistenten Keime besonders
konfrontiert und macht spezielle Maßnahmen erforderlich.
}
{}
{}
{}
{}


\index{Spitalskeime}\gT{Spitalskeime}: Keime mit "`angezüchteten"' Resistenzen

\gFazit{Unterscheide:
{\textquotedblleft}Wald-und-Wiesen{\textquotedblright}-Keime und
Spitalskeime} 

\gFazit{Auch im Rettungs- und Krankentransport kommen
gehäuft resistente Keime vor! Der "`Haltegriff in der
U-Bahn"' ist daher nicht mit dem "`Haltegriff im
KTW\marginlabel{\emph{KTW}: Krankentransportwagen}"'
vergleichbar.}

Der wohl bekannteste Problemkeim ist der \gT{MRSA}
(\gZ{Multiresistenter Staphylococcus aureus, siehe unten)}.
Das liegt schlicht daran dass er der erste war. In den letzten Jahren ist die Resistenzentwicklung aber
auch bei vielen anderen Keimen zu beobachten, \gZ{Beispiele
dafür sind die \gAbk{ESBL}{\textquotesingle}s und \gAbk{VRE}{\textquotesingle}s
sowie der VRSA (Extended Spectrum Beta-Laktamase-Bildner, Vancomycin-resistente
Enterokokken, Vancomycin-resistenter Staph. aureus). Auch bei der Tuberkulose
treten mittlerweile resistente Keime auf: \gE{2008} wurden die ersten Fälle
registriert bei denen der Keim auf mind. \gE{4} Antibiotika resistent
war\opt{NIVMED}{ (\gE{XDR-TBC}. Resistenzen gegen: Isoniazid, Rifampizin,
Fluochinolone, sowie zumindet eines der folgenden Antobiotika: Amikacin, Capreomycin, Kanamycin)}}\footnote{Quelle: AGES}.
% TODO LIT, HYG: MDR/XDR-TB Quelle bei AGES

\gEs{Das Problem der multiresistenten Keime wird immer größer.}

\gFazit{Je wirkungsloser die Antibiotika werden desto
wichtiger ist die Hygiene und die Vermeidung von
(Kreuz-)Infektionen.}

\begin{minipage}{13.67cm}
% [Warning: Draw object ignored]
% TODO BILD: MRSA Kreis (Grafik)
% \gCaptionof{figure}{Teufelskreis der
% resistenten Keime. Oft das Bindeglied: Der Rettungs- und
% Krankentransportdienst.
% \gSourcePicture{unbekannt} \gAuthNotes{Achtung! Zeichnung
% fehlt!}}


\end{minipage}

\subsection{MRSA -- Multi-Resistenter Staphylococcus
aureus} 

\gZ{Ursprünglich: Methicilin-Resistenter Staphylococcus Aureus}


\gLayout{28}{Staphylococcus aureus}
{
Der \emph{Staphylococcus aureus} ist ein
"`Allerweltskeim", bis zu 50\% der
Bevölkerung sind besiedelt. Vor allem ist die vordere
Nasenhöle besiedelt, aber auch die Haut (Achselhöle,
Analregion) und die Schleimhäute können ihn beherbergen.
Auch im Spital ist er häufig, aufgrund des hohen
Antibiotikaverbauchs ist er dort aber oft schon
Antibiotika-Resistent: Man spricht dann vom \gT{MRSA}.
}{
\begin{itemize}
    \item Staph aureus = "`Allerweltskeim"'
    \item aber resistente Form:
    \gT{MRSA}\footnote{Multi-resistenter Staphylococcus aureus}
    \begin{itemize}
        \item {\textquotedblleft}Spitalsinfektionen{\textquotedblright}
        \item schwer bis gar nicht behandelbar
        \item Wunden, die nicht abheilen wollen, Sepsis, Tod ...
    \end{itemize}
\end{itemize}
}
{./images/hygiene/11159_lores_VRSA.jpg}
{\gAuthNotesPicLegalOk{PD}{}Ein
resistenter Staph. aureus (\gAbk{VRSA}) im Elektronenmikroskop,
coloriert.}
{Public Health Image Library, ID\#: 11159, Public Domain.}



\gMassnahmenNG{Spezielle Maßnahmen}{MRSA-Transport}{\label{m:mrsa-transport}}{

\label{bkm:RefHeading40262819}

\begin{itemize}
    \item Informationen über MRSA-Trägerschaft
%     \marginlabel{gilt im Prinzip auch für die anderen
%     {\quotedblbase}Problemkeime{\textquotedblleft}}
    \item Art/Ort der MRSA-Besiedlung
    \item Patientenvorbereitung und Transport
    \item Wunden sind frisch verbunden und gut abgedeckt
    \item Bei Besiedlung der Atemwege trägt der Patient einen Mund-
    /Nasenschutz
    \item Vor dem Transport führt der Patient eine hygienische
    Händedesinfektion durch
    \item Allgemeine Hygienemaßnahmen
    \begin{itemize}
        \item Einmalhandschuhe und Schutzkittel.
        \item Nach Ablegen sofortige hygienische Händedesinfektion
        \begin{itemize}
            \item Nach Abschluss des Patiententransportes sind alle Materialien,
            Geräte, Instrumente und Flächen, welche direkten Kontakt mit dem
            Patienten hatten gemäß dem üblichen Hygieneplan zu
            desinfizieren.
        \end{itemize}
    \end{itemize}
    \item Kugelschreiber! 
    \item Alle waagerechten Oberflächen des Fahrzeuginnenraumes sind mit
    einem Flächendesinfektionsmittel in Form einer
    Scheuerwischdesinfektion desinfizierend zu reinigen.
    \item Abschließend führt das Einsatzpersonal eine hygienische
    Händedesinfektion durch.
    \item Nach Abschluss der vorangegangenen Hygiene- und
    Desinfektionsmaßnahmen ist das Fahrzeug wieder voll einsatzbereit.
\end{itemize}

Für die anderen Problemkeime \gZ{(VRSA, VRE, ESBL,
{\dots})} gelten sinngemäß die gleichen
Verfahrensweisen.
}



\section{Reinigung und hygienische Wiederaufbereitung}
\index{Desinfektion}

Man unterscheidet grundsätzlich zwischen \gE{Desinfetion} und 
\gE{Sterilisation}:

\begin{table}[H]
\begin{gWideParEnv}
\gCaptionof{table}{Gegenüberstellung Desinfektion -- Sterilisation.}

\begin{tabularx}{\textwidth}[H]{XX}
\gT{Desinfektion}
 &
\gT{Sterilisation}
\\\midrule


Ziel: \gEs{Keim\underline{reduktion}} ${\rightarrow}$ nicht
infektiös

Richtet sich gegen ausgewählte Keime

Aber: nicht keimfrei! nicht steril!

Chemische Substanzen

Einlegebad

 &

Abtötung aller Mikroorganismen -
{\textquotedblleft}\gEs{\underline{keimfrei}}{\textquotedblright}

Dampfsterilisation (Überdruck)

Heißluftsterilisation

\\\bottomrule
\end{tabularx}
\end{gWideParEnv}
\end{table}

\paragraph{Hygieneplan - Hygienebeauftragter}
\index{Hygieneplan}
\label{topic:hygieneplan}~

\gMarried{Der Hygieneplan regelt die konkreten Hygienemaßnahmen und Was Wann
Wie Womit (Wo) zu geschehen hat, siehe allgemeiner Aushang.}
{Der derzeitige Hygienebeauftragte ist

\gTakeNotesMedium}.


%%%%%%%%%%%%%%%%%%%%%%%%%%%%%%%%%%%%%%%%%%%%%%%%%%%%%%%%%%%
%%%%%%%%%%%%%%%%%%%%%%%%%%%%%%%%%%%%%%%%%%%%%%%%%%%%%%%%%%%
%%%%%%%%%%%%%%%%%%%%%%%%%%%%%%%%%%%%%%%%%%%%%%%%%%%%%%%%%%%
\subsection{Persönliche Hygiene}
\marginpar{\includegraphics[width=4cm]{images/AASdev20090228-img99.jpg}}
\label{topic:persoenliche-hygiene}

\begin{itemize}
    \item Schmuck, Uhren, Freundschaftsarmbänder sind perfekte
    Nährböden ${\rightarrow}$ Vor dem Dienst ablegen!
    \item Dienstkleidung regelmäßig waschen!
    \item Bei Kontamination wechseln!
    \item Verwendung von Plastikschürzen
    \item Bei Dienstantritt
    \begin{itemize}
        \item Hände waschen
        \item Nagelfalze reinigen
        \item Hygienische Händedesinfektion
        \index{Händedesinfektion}\index{Hygienische Händedesinfektion}
    \end{itemize}
\end{itemize}



%%%%%%%%%%%%%%%%%%%%%%%%%%%%%%%%%%%%%%%%%%%%%%%%%%%%%%%%%%%
%%%%%%%%%%%%%%%%%%%%%%%%%%%%%%%%%%%%%%%%%%%%%%%%%%%%%%%%%%%
%%%%%%%%%%%%%%%%%%%%%%%%%%%%%%%%%%%%%%%%%%%%%%%%%%%%%%%%%%%
\subsection{Flächen- und Wischdesinfektion}

\begin{itemize}
    \item Mit Flächendesinfektionsmittel \index{Flächendesinfektionsmittel}
    \item Einwirkzeit lt. Packung beachten!
    \item Fläche muss während der Zeit feucht sein
    \item \gEs{in Autos nur Flächendesinfektionsmittel
    ohne Alkohol!}
    \begin{itemize}
        \item Alkohol zerstört Acryl-Flächen! ${\rightarrow}$
        \index{Acryl-Des{\texttrademark}}Acryl-Des{\texttrademark}
    \end{itemize}
    \item nicht den Schmutz überall verteilen
    \begin{itemize}
        \item \gEs{Von außen in Richtung Schmutzquelle
        wischen}
        \item Einmaltücher wechseln
    \end{itemize}
\end{itemize}

\gCave{In den Fahrzeugen dürfen keine alkoholischen
Desinfektionsmittel verwendet werden, da diese die
Acryloberflächen zerstören können!}



%%%%%%%%%%%%%%%%%%%%%%%%%%%%%%%%%%%%%%%%%%%%%%%%%%%%%%%%%%%
%%%%%%%%%%%%%%%%%%%%%%%%%%%%%%%%%%%%%%%%%%%%%%%%%%%%%%%%%%%
%%%%%%%%%%%%%%%%%%%%%%%%%%%%%%%%%%%%%%%%%%%%%%%%%%%%%%%%%%%
\subsection{Gerätedesinfektion}
\gMarried{
\begin{itemize}
    \item Nach Kontamination:
    \begin{itemize}
        \item \gEs{Grobe Reinigung vor Ort}: Eintrocknen
        verhindern
        \item Wischdesinfektion
        \item Ggfs. Schläuche durchspülen
        \item Flächendesinfektionsmittel in Nierenschale / Joghurtbecher,
        saugen ...
    \end{itemize}
    \item Meldung an Leitstelle:
    \emph{{\textquotedblleft}Nicht EB, müssen
    putzen{\textquotedblright}}
    \item Kontaminierte Teile gegen Ersatzmaterial tauschen
    \item Zerlegen
    \item Nochmals spülen
    \item Wenn Desinfektionswanne leer ${\rightarrow}$ einlegen
    \begin{itemize}
        \item \gEs{Vollständig zerlegt} laut Anleitung
        \item \gEs{Luftblasenfrei!}
        \item Mit Niederhaltesieb unter die Oberfläche drücken
        \item Zeit auf Zettel notieren
        \begin{itemize}
            \item {\textquotedblleft}was, wann, von wo, von wem
            eingelegt{\textquotedblright}
        \end{itemize}
        \item sonst nur bereit legen, in Plastiksack
        \begin{itemize}
            \item {\textquotedblleft}was, von wo, von wem, warum{\textquotedblright}
        \end{itemize}
    \end{itemize}
\end{itemize}
}{
\includegraphics[width=\linewidth]{./images/AASdev20090228-img100.jpg}
    \gCaptionof{figure}{Desinfektionswannen, verschiedene
    Größen}
}

\begin{wrapfigure}{r}{4cm}

\end{wrapfigure}




% \clearpage
%%%%%%%%%%%%%%%%%%%%%%%%%%%%%%%%%%%%%%%%%%%%%%%%%%%%%%%%%%%
%%%%%%%%%%%%%%%%%%%%%%%%%%%%%%%%%%%%%%%%%%%%%%%%%%%%%%%%%%%
%%%%%%%%%%%%%%%%%%%%%%%%%%%%%%%%%%%%%%%%%%%%%%%%%%%%%%%%%%%
\subsection{Hygienische Händedesinfektion}

Die hygienische Händedesinfektion ist eine der wichtigsten Maßnahmen zur
Verhütung von Kreuzinfektionen. Sie ist eine Standardmaßnahme und wird unter
\prettyref{topic:hygienische-haendedesfinfektion} besprochen wird.


%%%%%%%%%%%%%%%%%%%%%%%%%%%%%%%%%%%%%%%%%%%%%%%%%%%%%%%%%%%
%%%%%%%%%%%%%%%%%%%%%%%%%%%%%%%%%%%%%%%%%%%%%%%%%%%%%%%%%%%
%%%%%%%%%%%%%%%%%%%%%%%%%%%%%%%%%%%%%%%%%%%%%%%%%%%%%%%%%%%
\subsection{Öffnen von sterilem Material}

Packung {\quotedblbase}aufschälen{\textquotedblleft} 

\begin{figure}[H]
\includegraphics[width=6.656cm,height=4.993cm]{images/AASdev20090228-img103.jpg}
\includegraphics[width=6.663cm,height=4.984cm]{images/AASdev20090228-img104.jpg}
\gCaptionof{figure}{Öffnen von sterilen Material unter verschiedenen
Bedingungen: Immer wird die Verpackung "`aufgeschält"'.
\gAuthNotesPicLegalNotOk{}{}}
\end{figure}



%%%%%%%%%%%%%%%%%%%%%%%%%%%%%%%%%%%%%%%%%%%%%%%%%%%%%%%%%%%
%%%%%%%%%%%%%%%%%%%%%%%%%%%%%%%%%%%%%%%%%%%%%%%%%%%%%%%%%%%
%%%%%%%%%%%%%%%%%%%%%%%%%%%%%%%%%%%%%%%%%%%%%%%%%%%%%%%%%%%
\subsection{Selbstschutz}

\begin{itemize}
    \item Immer Einmalhandschuhe beim Patientenkontakt tragen!
    \item Handschuhe nach jedem Patienten wechseln
    \item Nie mit schmutzigen Handschuhen in die Tasche greifen!
    \item \gEs{Ein Helfer hat Patientenkontakt, der andere
    reicht zu!}
    \item Achtung {\textquotedblleft}Kugelschreiber{\textquotedblright}
    \item Einmalschürzen bei Bedarf
    \begin{itemize}
        \item evt. Infektionsschutzanzüge, Maske
    \end{itemize}
    \item \gEs{Spitze Gegenstände (Nadeln, Lanzetten, ...
    sofort in die Kontamed-Box!}
    \begin{itemize}
        \item Ausnahme: ausdrückliche Arztanweisung (Blutzuckerbestimmung aus
        Nadel)
        \item \emph{{\textquotedblleft}Wer sticht, der
        entsorgt!{\textquotedblright}}
    \end{itemize}
\end{itemize}

\paragraph{Kontamed, Sharp, Gelbe Box ...}
Spitze Gegenstände (Nadeln, Lanzetten, ... sofort in Gelbe Box!

Ausnahme: ausdrückliche Arztanweisung (Blutzuckerbestimmung aus Nadel)

\gFazit{Wer sticht, der entsorgt!}



\begin{figure}[H]
\includegraphics[width=3.377cm,height=4.552cm]{images/AASdev20090228-img105.jpg}
\includegraphics[width=5.71cm,height=4.289cm]{images/AASdev20090228-img106.jpg}
\gCaptionof{figure}{Kontamed-Boxen, verschiedene Fabrikate.}
\end{figure}

\gFazit{NIE Nachstopfen.

NIE ausleeren.

Wenn die Box voll ist wird sie verschlossen und abgegeben.}



%%%%%%%%%%%%%%%%%%%%%%%%%%%%%%%%%%%%%%%%%%%%%%%%%%%%%%%%%%%
%%%%%%%%%%%%%%%%%%%%%%%%%%%%%%%%%%%%%%%%%%%%%%%%%%%%%%%%%%%
%%%%%%%%%%%%%%%%%%%%%%%%%%%%%%%%%%%%%%%%%%%%%%%%%%%%%%%%%%%
\subsection{Wundreinigung}

\gLayout{23}{}
{
\begin{itemize}
    \item Bei frischen, nicht infizierten Wunden:
    \begin{itemize}
        \item Reinigung von innen nach außen
        \item Ziel: keine Keime in Wunde bringen
    \end{itemize}
    \item Womit?
    \begin{itemize}
        \item Sterilen Lösungen oder Octenisept
        \item Sterile Tupfer
        \item KEINE Papierhandtücher oder unsterile
        Tupfer!
    \end{itemize}
    \item Danach steriler Verband!
    \item Bei offenen Frakturen ist Keimfreiheit be\-sonders wichtig !!!
    \begin{itemize}
        \item Gefahr einer lebenslangen Knochen\-marks\-ent\-zündung
        (Osteomyelitis)!
        \item \index{OpSite}\gT{OpSite}\texttrademark-Folie
        wenn vorhanden.
    \end{itemize}
\end{itemize}

Beachte die Grundregeln der Wundversorgung
\gSee{topic:wundversorgung-grundregeln} !
}{

}
{./images/noauth/AASdev20090228-img107.png}
{\gAuthNotesPicLegalNotOk{}{}Wundreinigung:
Bei einer frischen Wunde von innen nach außen.}
{}

%% Wird andernorts behandelt
% 
% \begin{figure}[H]
%     \includegraphics[width=4.219cm,height=4.219cm]{images/AASdev20090228-img108.jpg}
%     \caption{Anbringen
%     einer OpSite(TM)-Folie: Es ist wichtig die Folie zu spannen damit sie
%     sich nicht selbst verklebt.
%     \gAuthNotesPicLegalNotOk{}{} }
% \end{figure}


\clearpage
%%%%%%%%%%%%%%%%%%%%%%%%%%%%%%%%%%%%%%%%%%%%%%%%%%%%%%%%%%%
%%%%%%%%%%%%%%%%%%%%%%%%%%%%%%%%%%%%%%%%%%%%%%%%%%%%%%%%%%%
%%%%%%%%%%%%%%%%%%%%%%%%%%%%%%%%%%%%%%%%%%%%%%%%%%%%%%%%%%%
%%%%%%%%%%%%%%%%%%%%%%%%%%%%%%%%%%%%%%%%%%%%%%%%%%%%%%%%%%%
\section{Besondere Verhaltensweisen}

%%%%%%%%%%%%%%%%%%%%%%%%%%%%%%%%%%%%%%%%%%%%%%%%%%%%%%%%%%%
%%%%%%%%%%%%%%%%%%%%%%%%%%%%%%%%%%%%%%%%%%%%%%%%%%%%%%%%%%%
%%%%%%%%%%%%%%%%%%%%%%%%%%%%%%%%%%%%%%%%%%%%%%%%%%%%%%%%%%%
\subsection{Nadelstichverletzungen}\index{Nadelstichverletzungen}

\gMarginNo{\gZ{Basierend auf: Gemeinsame Erklärung der
Deutschen AIDS-Gesellschaft (\gAbk{DAIG}) und Osterreichischen
AIDS-Gesellschaft (\gAbk{OAG}). [Post-exposure prophylaxis of HIV
infection. German-Austrian recommendations, update
September 2007] \emph{Dtsch Med Wochenschr}, 2009, 134 Suppl
1, S16-S33}\cite{OeAG2007}}

\gMassnahmenNG{Spezielle Maßnahmen}{Nadelstichverletzung}{}{

\begin{itemize}
    \item Erstmaßnahmen durchführen: \nocite{OeAG2007}
    \begin{itemize}
        \item bei Stich-bzw.Schnittverletzungen:
        Blutung durch Kompression des umgebenden
        Gewebes inudzieren. Danach desinfizieren und einen
        in Desinfektionsmittel getränkten Tupfer 10 Minuten
        auf der Wunde belassen.
        \item bei Spritzern mit Blut oder Sekreten in Auge oder Mund:
        Schleimhaut mit Wasser u./o. Schleimhautdesinfektionsmittel spülen
    \end{itemize}
    \item Leitstelle informieren
    \item ${\rightarrow}$ \gEs{Umgehend in ein
    Unfallspital}\footnote{\ldots am besten in ein
    Unfallspital der \gAbk{AUVA}, da es sich um einen
    Arbeitsunfall handelt.}
    \item Patient: Zielspital informieren:
    \begin{itemize}
        \item Blut sichern für \gAbk{HIV}- und Hepatitis-Serologie
    \end{itemize}
    \item  \gC{Unfallmeldung}  (wichtig!), Blutabnahme
    \item ggf. \gAbk{HIV}-Postexpositionsprohylaxe (\gAbk{PEP})
    \begin{itemize}
        \item kurzfristige Therapie (4 Wo)
        \item Einzelfall muss immer individuell beurteilt werden
        (Nutzen/Risiko), starke Nebenwirkungen.
        \item Muss innerhalb von 2 Std. begonnen werden ${\rightarrow}$ daher
        \underline{SOFORT} in ein Unfallspital der \gAbk{AUVA}
    \end{itemize}
    \item Meldung an Betriebsleitung
    \item \gEs{Unfallmeldung (wichtig!) an die \gAbk{AUVA}}
\end{itemize}

\gCave{Unfallmeldung nicht vergessen!

Sofort in ein Unfallspital! }
}




\subsection{Schutzmaßnahmen bei infektiösen Patienten}

\subsubsection{Mund-/Nasenmasken}

Es gibt verschiedene Arten von Schutzmasken, welche sich bezüglich ihrer
Schutzwirkung wesentlich unterscheiden. Ein besonders wichtiges Kriterium ist 
die minimale Größe der Partikel welche eine Maske ausfiltern kann. Darüber
hinaus ist die Dichtigkeit zu beachten. 

\paragraph{}Die Einfachste Form ist die (chirurgische) Mund-Nasen-Maske
(\gAbk{MNM}, "`OP-Maske"'), welche für viele Situationen einen ausreichenden Schutz
darstellt\footnote{Beim Ausbruch von \gAbk{SARS} in Honkong in 2003 konnten
simple OP-Masken das Infektionsrisiko für das Personal auf 0 senken (das
infizierte Personal trug entweder keine Maske oder einfache Papiermasken). 
Seto W.H. et. al.: Lancet 2003
% TODO LIT: Seto W.H. et. al.: Lancet 2003
}.

\paragraph{FFP-Masken} Es gibt je nach Partikelgröße, die ausgefiltert wird,
verschiedene Klassen: FFP 1, FFP 2 und FFP 3. 
% TODO HYG: FFP-Masken
% TODO HYG: FFP-Masken Filtergrößen

\begin{table}[H]
\begin{gWideParEnv}
\begin{minipage}{\linewidth}
\gCaptionof{table}{Arten von Schutzmasken und Rückhaltevermögen.}

\renewcommand{\thefootnote}{\alph{footnote}}
\begin{tabulary}{\linewidth}{LLLLL}
\gTabHeading{Typ}
&
\gTabHeading{Mindestrückhaltevermögen des Filters}
\footnotemark[1]\footnotemark[2]
% TODO LIT: Dreller S. et.al.: RdL 66 (März 2006)
&
\gTabHeading{Max. zul.
Gesamtleckage}\footnotemark[3] &
\gTabHeading{Min. Partikelgröße}
&
\gTabHeading{Anwendung (Beispiele)}
\\\midrule\addlinespace
\gTabSub{Chir. Maske} &
95\,\%\footnotemark[4] & k.A. & & Basisschutz; \gAbk{MRSA}, \ldots
\\\addlinespace \gTabSub{FFP 1} & 80\,\% & 22\,\% & &
\\\addlinespace
\gTabSub{FFP 2} & 94\,\% & 8\,\% & & Influenza
\\\addlinespace
\gTabSub{FFP 3} & 99\,\% & 9\,\% & & \gAbk{SARS}, offene Tuberkulose
\\\addlinespace\bottomrule
\end{tabulary}
\footnotetext[1]{für \ce{NaCl}-Prüfaerosol}
\footnotetext[2]{Dreller S. et.al.: RdL 66 (März 2006)}
\footnotetext[3]{gem. DIN EN 149\label{fn:hyg-sm-dinen149}}
\footnotetext[4]{für Staphylococcus
aureus\label{fn:hyg-sm-stapha}}
\renewcommand{\thefootnote}{\arabic{footnote}}
\end{minipage}
\end{gWideParEnv}\end{table} 
%%%%%%%%%%%%%%%%%%%%%%%%%%%%%%%%%%%%%%%%%%%%%%%%%%%%%%%%%%%
%%%%%%%%%%%%%%%%%%%%%%%%%%%%%%%%%%%%%%%%%%%%%%%%%%%%%%%%%%%
%%%%%%%%%%%%%%%%%%%%%%%%%%%%%%%%%%%%%%%%%%%%%%%%%%%%%%%%%%%
\subsection{MRSA und andere Resistente Keime}
\index{MRSA}

Die Maßnahmen beim Transport von Patienten mit resistenten Keimen
werden unter \prettyref{m:mrsa-transport} besprochen.





\nocite{KayserMikrobiologie10}
\nocite{HahnMikrobio6}
\nocite{BLIM3}
\nocite{OHCSP7}
\nocite{OeAG2007}


\clearpage
\gChapterBibliography
